\documentclass[12pt]{article}

\usepackage[margin=1in]{geometry}
\IfFileExists{newtxtext.sty}{
    \usepackage{newtxtext,newtxmath}
}{
    \usepackage{mathptmx}
}
\usepackage{setspace}
\usepackage{amsmath}
\usepackage{graphicx}
\usepackage{booktabs}
\usepackage{cite}
\usepackage[hidelinks]{hyperref}

\doublespacing
\setlength{\parindent}{0.5in}
\setlength{\parskip}{0pt}

\begin{document}

\begin{center}
    {\large\bfseries Modeling MOSFET Gate-Oxide Tunneling with a One-Dimensional Slanted Potential Barrier}
\end{center}

\section{Introduction}
Metal–oxide–semiconductor field-effect transistors (MOSFETs) are widely used as electronic switches because the voltage applied to the gate controls the conductivity of the channel. Oxide layer is placed between the gate and the silicon semiconductor to prevent electron transfer among them. Reducing the oxide thickness increases the oxide capacitance, allowing the gate voltage to control the channel charge more effectively.. However, in the micro-world with unit of \(10^{-9} \text{ m}\), thinner oxide also means a narrower potential barrier, thus more electrons engage in tunneling and increase gate leakage current. In classical model, an electron with insufficient energy (greater than work function) cannot cross the oxide barrier; quantum mechanically, its wavefunction extends out of the barrier, producing a nonzero probability of transmission. This paper will establish a model of leakage using Schrodinger's equation, test the model with experimental data, and discuss how this leakage affect the efficiency, thermal stability, and reliability.
\section{Theoretical Background}

\subsection{One-Dimensional Slanted Potential Barrier}

\subsection{WKB Approximation}

\section{Model Construction}

\subsection{Barrier Definition and Assumptions}

\subsection{Python Implementation}

\section{Experimental Data and Processing}

\section{Results}

\section{Discussion}

\subsection{Comparison with Experimental Gate Leakage}

\subsection{Effect of Leakage-Related Heating on MOSFET Operation}

\section{Evaluation}

\section{Conclusion}

\appendix
\section{Python Code}

\section{Additional Data and Calculations}

\newpage
\IfFileExists{IEEEtran.bst}{\bibliographystyle{IEEEtran}}{\bibliographystyle{unsrt}}
\bibliography{references}

\end{document}
